% Abstract (English then Hebrew). Per BGU \S 10.7.2.3.
% PRD V2 \S 7.2 R1, R5, Q4, Q5, FX-16. The 80 \% CIFAR-10 acceptance
% criterion is NOT mentioned here (R1) -- it appears once, in the
% Conclusions section, against the as-measured numbers. Cross-model
% magnitudes are the as-measured values (PRD \S 12 Q2 escape clause).
\section*{Abstract}
\addcontentsline{toc}{section}{Abstract}

\noindent\textbf{Background.}\quad
Terahertz (THz) imaging sensors produce frames that are simultaneously low
resolution, low signal-to-noise, and colorless. ImageNet-pretrained
classifiers, which rely on fine textures and saturated chroma, lose
accuracy quickly under those conditions, but the magnitude and the
mechanism of the loss across modern deep-learning architectures are not
documented quantitatively.

\noindent\textbf{Aim.}\quad
Compare the robustness of three representative deep architectures (a
residual CNN, a densely connected CNN, and a foveal vision transformer)
under controlled synthetic THz-like degradation, decompose the loss into
per-axis contributions, and quantify how much of it can be recovered by
training-time regularization.

\noindent\textbf{Methods.}\quad
The 402-cell campaign trains ResNet50, DenseNet121, and TransNeXt-tiny
on CIFAR-10 and MNIST under a five-level degradation curve (saturation,
downsample, blur, additive noise, salt and pepper). Seven phases isolate
the effects: A clean baseline, B1 combined, C single-axis, D
regularization recovery, B2 and B2-nr THz-protocol simplification, and
C2 single-axis under the THz protocol. A 48-replicate multi-seed audit
on the 24 L3 headline cells bounds reproducibility. Every degradation
pipeline is anchored to a deterministic 256-image PSNR / SSIM probe so
the cross-phase comparison sits on a physical scale.

\noindent\textbf{Findings.}\quad
Under Phase B1 on CIFAR-10, TransNeXt-tiny holds the highest accuracy
margin over the two CNNs at the mild and moderate end of the curve
($+10.1\,$\Mpp{} at L1, $+10.4\,$\Mpp{} at L2 vs the CNN mean), and the
margin narrows to under $1.5\,$\Mpp{} from L3 onward. Resolution
dominates the per-axis attribution by a factor of approximately four
(mean L1 to L5 effect on CIFAR-10: $-27.45\,$\Mpp{} for resolution vs
$-6.58$ for blur and $-5.13$ for salt and pepper). Regularization (T3 in
Phase D) recovers approximately $0.77\,$\Mpp{} of the $\sim 12.5\,$\Mpp{}
Phase B1 accuracy collapse. The multi-seed audit confirms
$\sigma < 1.5\,$\Mpp{} on every L3 headline cell.

\noindent\textbf{Conclusions.}\quad
The accuracy collapse is best modeled as an information bottleneck, not
as an overfit. The foveal transformer holds the highest robustness
margin at mild and moderate degradation; all three architectures
converge into the $[19, 21]\,$\Mpp{} band at L5. Future THz deployments
should invest in resolution restoration before denoising or deblurring.

\noindent\textbf{Keywords.}\quad
object classification, deep learning, low resolution, THz imaging,
robustness, PSNR, SSIM, calibration, transformer, residual network,
dense network.

\bigskip
\hrule
\bigskip

% ====================================================================
% Hebrew abstract (PRD V2 \S 7.2 R5, Q5, FX-16). Inlined here so the
% DOCX deliverable carries the Hebrew text in the body, not in an
% appendix. Rendered by XeLaTeX + polyglossia; pdflatex falls back to
% a placeholder pointing at the XeLaTeX render. The DOCX build
% (Pandoc + reference template) walks the source text directly and
% renders Hebrew correctly regardless of the engine used here.
% ====================================================================
\ifdefined\XeTeXversion
  \begin{hebrew}
  % Hebrew abstract body — six labelled paragraphs mirroring the English
% structure. PRD V2 §7.2 Q5, FX-16. Operator reviews before D5 QA.

\noindent\textbf{רקע.}\quad
חיישני הדמיה בתחום הטרה-הרץ (THz) מפיקים תמונות בעלות רזולוציה נמוכה,
יחס אות-לרעש נמוך והעדר צבע. מסווגי \emph{ImageNet} שאומנו מראש,
המסתמכים על מרקמים עדינים ועל צבע רווי, מאבדים דיוק במהירות תחת תנאים
אלה, אך עוצמת הירידה והמנגנון העומד בבסיסה אינם מתועדים כמותית עבור
שלוש משפחות הארכיטקטורות המודרניות.

\noindent\textbf{מטרה.}\quad
להשוות את עמידותן של שלוש ארכיטקטורות-עומק מייצגות (רשת CNN שאריתית,
רשת CNN צפופה, וטרנספורמר ראייה פוֹביאַלי) תחת השפלה סינתטית מבוקרת
דמוית-THz, לפרק את ההפסד לתרומות לפי ציר, ולכמת כמה ממנו ניתן לשחזר על
ידי רגולריזציית-אימון.

\noindent\textbf{שיטות.}\quad
המערכה כוללת 402 תאי אימון של ResNet50, DenseNet121 ו-TransNeXt-tiny על
CIFAR-10 ו-MNIST, תחת עקומת השפלה בת חמש רמות (רוויה, הקטנת רזולוציה,
טשטוש, רעש גאוסי, מלח-ופלפל). שבעה שלבים מבודדים את האפקטים: A קו בסיס
נקי, B1 השפלה משולבת, C ציר יחיד, D שחזור באמצעות רגולריזציה, B2 ו-B2-nr
לפישוט פרוטוקול ה-THz, ו-C2 ציר יחיד תחת אותו פרוטוקול. בדיקת רבייה
על 48 חזרות עם זרעי-אקראיות נוספים שולחת חסם על שונות-בין-ריצות בכל אחד
מ-24 תאי הכותרת ברמה L3. כל צינור השפלה מעוגן לבדיקת PSNR / SSIM
דטרמיניסטית על 256 תמונות, כך שההשוואה בין השלבים מתבצעת על סקאלה
פיזיקלית.

\noindent\textbf{ממצאים.}\quad
תחת שלב B1 ב-CIFAR-10, TransNeXt-tiny שומר על המרווח הרחב ביותר מעל שתי
רשתות ה-CNN בקצה המתון של העקומה (כ-$10.1\,$\Mpp{} ברמה L1, כ-$10.4\,$\Mpp{}
ברמה L2 מול הממוצע של ה-CNN-ים), ומרווח זה מצטמצם לפחות מ-$1.5\,$\Mpp{}
החל מ-L3. ציר הרזולוציה דומיננטי בייחוס לפי ציר בקירוב פי ארבעה
(אפקט ממוצע L1 עד L5 על CIFAR-10: $-27.45\,$\Mpp{} עבור רזולוציה לעומת
$-6.58$ עבור טשטוש ו-$-5.13$ עבור מלח-ופלפל). רגולריזציה (T3 בשלב D)
משחזרת בקירוב $0.77\,$\Mpp{} מתוך קריסת הדיוק של $\sim 12.5\,$\Mpp{} בשלב
B1. בדיקת הרבייה מאשרת $\sigma < 1.5\,$\Mpp{} בכל תא כותרת ברמה L3.

\noindent\textbf{מסקנות.}\quad
קריסת הדיוק מתוארת היטב כצוואר-בקבוק של מידע ולא כהתאמת-יתר. הטרנספורמר
הפוֹביאַלי שומר על מרווח העמידות הגבוה ביותר בקצה המתון של עקומת ההשפלה;
שלוש הארכיטקטורות מתכנסות לרצועת $[19, 21]\,$\Mpp{} ברמה L5. פריסות THz
עתידיות צריכות להשקיע בשחזור רזולוציה לפני שלבי הסרת רעש או הסרת טשטוש.

\noindent\textbf{מילות מפתח.}\quad
סיווג עצמים, למידה עמוקה, רזולוציה נמוכה, הדמיית THz, עמידות, PSNR, SSIM,
כיול, טרנספורמר, רשת שאריתית, רשת צפופה.

  \end{hebrew}
\else
  \noindent\emph{\small Hebrew abstract: see XeLaTeX render or
  \texttt{sections/hebrew\_abstract\_body.tex} for the inline source.}
\fi
